% !TeX root = ../example.tex
% Four subsections retain the two-column navigation example.
\section{导航与渐进显示}
\sectioncoverpage[assets/tsinghua_door1]{本节演示逐步显示、内容替换、导航跳转与换页效果。}
\note{本节仍保留四个小节的双栏排列。前两页分别演示累积显示与固定区域内的内容替换，随后说明跳转和导出。}

\subsection{逐步显示}
\begin{frame}[fragile,label=demo-reveal]{按讲述顺序呈现要点}
  \transfade<2-|handout:0>[duration=0.2]
  \begin{enumerate}
    \item<1-> \textbf{问题}：先明确这一页要回答什么。
    \item<2-> \textbf{证据}：再给出支持判断的材料。
    \item<3-> \textbf{结论}：最后收束到一个核心信息。
  \end{enumerate}
  \medskip
  \small 用 \verb|\item<1->|、\verb|\item<2->|、\verb|\item<3->| 保留已显示内容。
  \note<1>{逐步显示，第 1 步。先解释问题；后两条尚未出现，但位置已保留。此帧的第一阶段不设置换页效果。}
  \note<2>{逐步显示，第 2 步。增加证据，已显示的问题保持原位。PDF 的淡入设置会导出为 PPTX 原生淡入换页，仍由点击推进。}
  \note<3>{逐步显示，第 3 步。增加结论并保留前两条；此阶段是完整状态。每一步使用各自的 TeX 备注。}
\end{frame}

\subsection{内容替换}
\begin{frame}[fragile,label=demo-replace]{在固定区域内替换内容}
  \transdissolve<2-|handout:0>[duration=0.4]
  \begin{overlayarea}{\textwidth}{1.8cm}
    \only<1|handout:0>{%
      \textbf{输入：自由草稿}\par\medskip
      问题、证据与讲述提示还没有分层。
    }
    \only<2|handout:1>{%
      \textbf{输出：详细大纲}\par\medskip
      按 section/subsection 组织要点、证据与备注安排。
    }
  \end{overlayarea}
  \small 用 \verb|overlayarea| 保留区域；用 \verb|\only| 切换内容。
  \note<1>{内容替换，第 1 步。这里只显示输入状态。固定高度的 overlayarea 保留下一阶段的位置，避免页脚或后续文字随内容长度移动。}
  \note<2>{内容替换，第 2 步。同一位置替换为输出状态，PPTX 保留溶解换页。handout 只保留此最终状态；草稿内容的实质扩充仍在详细大纲阶段完成。}
\end{frame}

\subsection{导航跳转}
\begin{frame}{导航与帧标签}
  \begin{description}[小节跳转]
    \item[章节跳转] 页脚上层定位到对应章节。
    \item[小节跳转] 页脚下层定位到本节主题。
    \item[帧内标签] \hyperlink{demo-reveal}{返回逐步显示}\quad
      \hyperlink{demo-notes}{查看备注写法}
  \end{description}
\end{frame}
\note{点击两个帧标签和页脚导航，检查 PDF 与 PPTX 的目标一致。返回逐步显示时定位到第一阶段；备注写法在下一节。四个小节仍交替排列在章节页左右两栏。}

\subsection{输出规则}
\begin{frame}{展示阶段与换页}
  \begin{description}[PPTX 换页]
    \item[PDF 阶段] 每个 overlay 对应一个可独立查看的页面。
    \item[PPTX 换页] 每个阶段对应一张幻灯片，保留淡入或溶解设置。
    \item[讲义版本] \texttt{handout} 合并展示阶段，替换区域保留最终内容。
  \end{description}
  \medskip
  默认点击推进，按讲解节奏控制展示。
\end{frame}
\note{这两组示例使用 transfade 和 transdissolve，并以 handout:0 排除讲义版换页效果。PPTX 将每个 PDF 阶段作为单独幻灯片，并保留支持的原生换页设置；没有把图像中的文字变成独立动画对象。这里的页脚仍按 Beamer 逻辑帧编号，同一帧的多个阶段共用帧号。}
